\documentclass[../Main.tex]{subfiles}
\begin{document}
  \chapter{2021-2022学年微积分（一）（下）期末考试}

  \section{单项选择题(每小题 3 分，共 18 分)}
  \begin{enumerate}
    \item 微分方程 \( y'' + y = x^2 + 1 + \sin x \) 的待定特解形式可设为【\quad】.

      \begin{tasks}[label=\textbf{\Alph*.}, label-width=1.5em, column-sep=1em](2)
        \task \( y^* = ax^2 + bx + c + x(A\sin x + B\cos x) \)
        \task \( y^* = ax^2 + bx + c + A\sin x \)
        \task \( y^* = x(ax^2 + bx + c + A\sin x + B\cos x) \)
        \task \( y^* = ax^2 + bx + c + A\cos x \)
      \end{tasks}

    \item 在下列极限结果中，正确的是【\quad】.

      \begin{tasks}[label=\textbf{\Alph*.}, label-width=1.5em, column-sep=1em](1)
        \task \( \lim\limits_{\substack{x \to 0 y \to 0}} \dfrac{xy}{x^2 + y^2} = 0 \)
        \task \( \lim\limits_{\substack{x \to 0 y \to 0}} \dfrac{x^2 y}{x^2 + y^2} = 0 \)
        \task \( \lim\limits_{\substack{x \to 0 y \to 0}} \dfrac{xy}{x + y} = 0 \)
        \task \( \lim\limits_{\substack{x \to 0 y \to 0}} \dfrac{x^2 y}{x + y} = 0 \)
      \end{tasks}

    \item 设 \( f(x) \) 为连续函数，\( F(t) = \displaystyle \int_1^t \mathrm{d}y \displaystyle \int_y^t f(x) \mathrm{d}x \)，则 \( F'(2) \) 等于【\quad】.

      \begin{tasks}[label=\textbf{\Alph*.}, label-width=1.5em, column-sep=1em](2)
        \task \( 2f(2) \)
        \task \( f(2) \)
        \task \( -f(2) \)
        \task \( 0 \)
      \end{tasks}

    \item 设曲线 \( L: f(x, y) = 1 \)（\( f(x, y) \) 具有一阶连续偏导数），过第 II 象限内的点 \( M \) 和第 IV 象限内的点 \( N \)，\( T \) 为 \( L \) 上从点 \( M \) 到点 \( N \) 的一段弧，则下列积分小于零的是【\quad】.

      \begin{tasks}[label=\textbf{\Alph*.}, label-width=1.5em, column-sep=1em](1)
        \task \( \displaystyle \int_T f(x, y) \mathrm{d}s \)
        \task \( \displaystyle \int_T f(x, y) \mathrm{d}x \)
        \task \( \displaystyle \int_T f(x, y) \mathrm{d}y \)
        \task \( \displaystyle \int_T f_x'(x, y) \mathrm{d}x + f_y'(x, y) \mathrm{d}y \)
      \end{tasks}

    \item 设 \( L \) 为空间曲线 \( \begin{cases} x^2 + y^2 + z^2 = a^2 \\ x + y + z = 0 \end{cases} \)，则 \( \displaystyle \oint_L x^2 \mathrm{d}s = \)【\quad】.

      \begin{tasks}[label=\textbf{\Alph*.}, label-width=1.5em, column-sep=1em](2)
        \task \( 0 \)
        \task \( 2\pi a^3 \)
        \task \( \dfrac{1}{3}\pi a^2 \)
        \task \( \dfrac{2}{3}\pi a^3 \)
      \end{tasks}

    \item 设两个数列 \( \{a_n\}, \{b_n\} \)，若 \( \lim\limits_{n \to \infty} a_n = 0 \)，则【\quad】.

      \begin{tasks}[label=\textbf{\Alph*.}, label-width=1.5em, column-sep=1em](1)
        \task 当 \( \sum \limits_{n=1}^{\infty} b_n \) 收敛时，\( \sum \limits_{n=1}^{\infty} a_n b_n \) 收敛
        \task 当 \( \sum \limits_{n=1}^{\infty} b_n \) 发散时，\( \sum \limits_{n=1}^{\infty} a_n b_n \) 发散
        \task 当 \( \sum \limits_{n=1}^{\infty} |b_n| \) 收敛时，\( \sum \limits_{n=1}^{\infty} a_n^2 b_n^2 \) 收敛
        \task 当 \( \sum \limits_{n=1}^{\infty} |b_n| \) 发散时，\( \sum \limits_{n=1}^{\infty} a_n^2 b_n^2 \) 发散
      \end{tasks}
  \end{enumerate}

  \section{填空题(每小题 4 分，共 16 分)}
  \begin{enumerate}
    \setcounter{enumi}{6}
    \item 经过点 \( A(1,-2,3) \) 并且包含 \( x \) 轴的平面方程为 \(\underline{\hspace{5em}}\).
      \vspace{1em}

    \item 设矢量场 \( \boldsymbol{A} = x^2 \boldsymbol{i} + y z \boldsymbol{j} + z x \boldsymbol{k} \)，则 \(  \text{rot}\boldsymbol{A} \Big|_{(1,1,2)} = \underline{\hspace{5em}}\).
      \vspace{1em}

    \item \( u = x e^{y}z^3 \) 在点 \( (1,1,1) \) 处的全微分 \( \mathrm{d}u = \underline{\hspace{5em}}\).
      \vspace{1em}

    \item 若将函数 \( f(x) = \pi - x \)（\( 0 \leq x \leq \pi \)）展开成正弦级数 \( \sum \limits_{n=1}^{\infty} b_n \sin nx \)，则系数 \( b_4 \) 的值为 \(\underline{\hspace{5em}}\).
      \vspace{1em}
  \end{enumerate}

  \section{基本计算题(每小题 7 分，共 42 分)}

  \begin{enumerate}
    \setcounter{enumi}{10}
    \item 设方程组 \( \begin{cases} u + v = x \\ u^2 + v^2 = y \end{cases} \) 确定隐函数 \( u = u(x,y), v = v(x,y) \)，且 \( u \neq v \)，求 \( \dfrac{\partial u}{\partial x} \)，\( \dfrac{\partial v}{\partial x} \).
      \vspace{9em}

    \item 设 \( \boldsymbol{n} \) 是曲面 \( S: z = x^2 + y^2 \) 在点 \( P_0(1,1,2) \) 处指向上侧的法矢量，求函数 \( u = xz^3 - 3yz \) 在点 \( P_0 \) 处沿方向 \( \boldsymbol{n} \) 的方向导数 \( \dfrac{\partial u}{\partial \boldsymbol{n}} \).
      \vspace{9em}

    \item 计算 \( I = \displaystyle \iint \limits_D \max\{xy, 1\} \mathrm{d}x\,\mathrm{d}y \)，其中 \( D \) 是正方形区域 \( 0 \leq x \leq 2, 0 \leq y \leq 2 \).
      \vspace{9em}

    \item 计算 \( I = \displaystyle \iint \limits_{\Sigma} (xy + yz + zx) \mathrm{d}S \)，其中 \( \Sigma \) 为 \( z = \sqrt{x^2 + y^2} \) 被 \( x^2 + y^2 = 2ax \) 所截得的有限部分.
      \vspace{9em}

    \item 设 \( S \) 是球面 \( x^2 + y^2 + z^2 = R^2 \)（\( R > 0 \)）的外侧，求 \( I = \displaystyle \iint \limits_S x^3 \mathrm{d}y \mathrm{d}z + y^3 \mathrm{d}z \mathrm{d}x + z \mathrm{d}x \mathrm{d}y \).
      \vspace{9em}

    \item 确定幂级数 \( \sum \limits_{n=1}^{\infty} \dfrac{n}{n+1} x^n \) 的收敛域并求其和函数 \( S(x) \).
      \vspace{9em}
  \end{enumerate}

  \section{综合题(每小题 7 分，共 14 分)}
  \begin{enumerate}
    \setcounter{enumi}{16}
    \item 若二阶常系数线性齐次微分方程 \( y'' + ay' + by = 0 \) 的通解为 \( y = (C_1 + C_2 x)e^x \)，求非齐次方程 \( y'' + ay' + by = x \) 满足条件 \( y(0) = 2 \)，\( y'(0) = 0 \) 的解.
      \vspace{10em}

    \item 求 \( f(x, y) = x^2 + 2y^2 - x^2 y^2 \) 在区域 \( D = \{(x, y) \mid x^2 + y^2 \leq 4, y \geq 0\} \) 上的最大值和最小值.
      \vspace{10em}
  \end{enumerate}

  \section{证明题(每小题 5 分，共 10 分)}

  \begin{enumerate}
    \setcounter{enumi}{18}
    \item 证明级数 \( \sum \limits_{n=1}^{\infty} \dfrac{(-1)^n}{n - (-1)^n} \) 是条件收敛的.
      \vspace{10em}

    \item 设函数 \( f(x) \) 在 \( (-\infty, +\infty) \) 内具有一阶连续导数，\( L \) 是上半平面 \( (y > 0) \) 内的有向分段光滑曲线，其始点为 \( (a, b) \)，终点为 \( (c, d) \).记 \( I = \displaystyle \int_L \dfrac{1}{y}[1 + y^2 f(xy)]\mathrm{d}x + \dfrac{x}{y^2}[y^2 f(xy) - 1]\mathrm{d}y \)，(1) 证明曲线积分 \( I \) 与路径 \( L \) 无关；(2) 当 \( ab = cd \) 时，求 \( I \) 的值.
      \vspace{10em}
  \end{enumerate}
\chapter{2021-2022学年微积分（一）（下）期末考试参考答案}
  \section{单项选择题(每小题 3 分，共 18 分)}
  \begin{enumerate}
    \item \textbf{Solution}. A.

      多项式项 \(x^2+1\) 的特解可设为二次多项式，而 \(\sin x\) 与齐次方程
      \[
        y''+y=0
      \]
      的解共振，因此三角项须乘以 \(x\). 所以待定特解应设为
      \[
        y^*=ax^2+bx+c+x(A\sin x+B\cos x).
      \]

    \item \textbf{Solution}. B.

      由
      \[
        \left|\frac{x^2y}{x^2+y^2}\right|\le |y|
      \]
      可知当 \((x,y)\to(0,0)\) 时该极限为 \(0\). 其余选项均可取不同路径说明极限不存在或不一定为 \(0\).

    \item \textbf{Solution}. B.

      先交换积分顺序：
      \[
        F(t)=\int_1^t\,\mathrm{d}y\int_y^t f(x)\,\mathrm{d}x=\int_1^t(x-1)f(x)\,\mathrm{d}x.
      \]
      因而
      \[
        F'(t)=(t-1)f(t),
      \]
      所以
      \[
        F'(2)=f(2).
      \]

    \item \textbf{Solution}. C.

      因为在曲线 \(L:f(x,y)=1\) 上恒有 \(f(x,y)=1\)，故
      \[
        \int_T f(x,y)\,\mathrm{d}y=\int_T \mathrm{d}y=y_N-y_M<0,
      \]
      其中 \(M\) 在第二象限、\(N\) 在第四象限，所以 \(y_N<y_M\). 其余几个积分不小于零或恒为零.

    \item \textbf{Solution}. D.

      交线 \(L\) 是球面与过原点平面的交圆，半径为 \(a\). 由对称性
      \[
        \oint_L x^2\,\mathrm{d}s=\oint_L y^2\,\mathrm{d}s=\oint_L z^2\,\mathrm{d}s.
      \]
      又
      \[
        \oint_L(x^2+y^2+z^2)\,\mathrm{d}s=a^2\oint_L\,\mathrm{d}s=a^2\cdot 2\pi a=2\pi a^3.
      \]
      三者相等，所以
      \[
        \oint_L x^2\,\mathrm{d}s=\frac{2\pi a^3}{3}.
      \]

    \item \textbf{Solution}. C.

      由 \(a_n\to 0\) 可知充分大时 \(|a_n|\le 1\)，于是
      \[
        a_n^2b_n^2\le |b_n|.
      \]
      若 \(\sum |b_n|\) 收敛，则由比较判别法 \(\sum a_n^2b_n^2\) 收敛. 其余命题都不成立.
  \end{enumerate}

  \section{填空题(每小题 4 分，共 16 分)}
  \begin{enumerate}
    \setcounter{enumi}{6}
    \item \textbf{Solution}. $3y+2z=0$.

      由于所求平面包含 \(x\) 轴，所以它含有方向向量 \((1,0,0)\). 又它经过点 \(A(1,-2,3)\)，故还含有向量 \((1,-2,3)\). 两向量叉积可取法向量
      \[
        (1,0,0)\times(1,-2,3)=(0,-3,-2),
      \]
      所以平面方程为
      \[
        3y+2z=0.
      \]

    \item \textbf{Solution}. $-\boldsymbol{i}-2\boldsymbol{j}$.

      \[
        \boldsymbol{A}=(x^2,yz,zx),
      \]
      因而
      \[
        \operatorname{rot}\boldsymbol{A}
        =\left(\frac{\partial(zx)}{\partial y}-\frac{\partial(yz)}{\partial z},
        \frac{\partial(x^2)}{\partial z}-\frac{\partial(zx)}{\partial x},
        \frac{\partial(yz)}{\partial x}-\frac{\partial(x^2)}{\partial y}\right)
        =(-y,-z,0).
      \]
      在 \((1,1,2)\) 处取值即得
      \[
        -\boldsymbol{i}-2\boldsymbol{j}.
      \]

    \item \textbf{Solution}. $e\,\mathrm{d}x+e\,\mathrm{d}y+3e\,\mathrm{d}z$.

      \[
        u=xe^yz^3,\quad u_x=e^yz^3,\quad u_y=xe^yz^3,\quad u_z=3xe^yz^2.
      \]
      在 \((1,1,1)\) 处有 \(u_x=u_y=e,\ u_z=3e\)，故
      \[
        \mathrm{d}u=e\,\mathrm{d}x+e\,\mathrm{d}y+3e\,\mathrm{d}z.
      \]

    \item \textbf{Solution}. $\dfrac12$.

      由
      \[
        b_n=\frac{2}{\pi}\int_0^\pi(\pi-x)\sin nx\,\mathrm{d}x,
      \]
      当 \(n=4\) 时分部积分可得
      \[
        b_4=\frac12.
      \]
  \end{enumerate}

  \section{基本计算题(每小题 7 分，共 42 分)}
  \begin{enumerate}
    \setcounter{enumi}{10}
    \item \textbf{Solution}. 对方程组
      \[
        u+v=x,\qquad u^2+v^2=y
      \]
      关于 \(x\) 求偏导，得
      \[
        u_x+v_x=1,\qquad 2uu_x+2vv_x=0.
      \]
      解得
      \[
        u_x=\frac{v}{v-u},\qquad v_x=-\frac{u}{v-u}.
      \]

    \item \textbf{Solution}. 曲面写成
      \[
        z-x^2-y^2=0,
      \]
      所以上侧法向量可取
      \[
        \boldsymbol{n}=(-2x,-2y,1).
      \]
      在 \(P_0(1,1,2)\) 处，
      \[
        \boldsymbol{n}=(-2,-2,1),\qquad \boldsymbol{n}_0=\frac13(-2,-2,1).
      \]
      又
      \[
        u=xz^3-3yz,\qquad \nabla u=(z^3,-3z,3xz^2-3y),
      \]
      所以
      \[
        \nabla u(P_0)=(8,-6,9).
      \]
      因而方向导数为
      \[
        \frac{\partial u}{\partial \boldsymbol{n}}\Big|_{P_0}
        =\nabla u(P_0)\cdot \boldsymbol{n}_0
        =\frac{-16+12+9}{3}
        =\frac53.
      \]

    \item \textbf{Solution}. 按曲线 \(xy=1\) 将区域分开. 当 \(xy<1\) 时 \(\max\{xy,1\}=1\)，当 \(xy\ge 1\) 时 \(\max\{xy,1\}=xy\). 故
      \[
        I=\iint_{D_1}xy\,\mathrm{d}x\,\mathrm{d}y+\iint_{D_2}1\,\mathrm{d}x\,\mathrm{d}y.
      \]
      逐段积分得
      \[
        I=\int_{1/2}^2x\,\mathrm{d}x\int_{1/x}^2y\,\mathrm{d}y
        +\int_0^{1/2}\,\mathrm{d}y\int_0^2\,\mathrm{d}x
        +\int_{1/2}^2\,\mathrm{d}x\int_0^{1/x}\,\mathrm{d}y
        =\frac{19}{4}+\ln2.
      \]

    \item \textbf{Solution}. 曲面 \(\Sigma:z=\sqrt{x^2+y^2}\) 关于 \(xOz\) 面对称，而 \(xy\) 与 \(yz\) 都关于 \(y\) 为奇函数，因此
      \[
        \iint_\Sigma xy\,\mathrm{d}S=0,\qquad \iint_\Sigma yz\,\mathrm{d}S=0.
      \]
      从而
      \[
        I=\iint_\Sigma zx\,\mathrm{d}S.
      \]
      投影到 \(xOy\) 面，因
      \[
        z_x=\frac{x}{\sqrt{x^2+y^2}},\qquad z_y=\frac{y}{\sqrt{x^2+y^2}},
      \]
      所以
      \[
        \mathrm{d}S=\sqrt{1+z_x^2+z_y^2}\,\mathrm{d}x\,\mathrm{d}y=\sqrt2\,\mathrm{d}x\,\mathrm{d}y.
      \]
      用极坐标 \(x=r\cos\theta,\ y=r\sin\theta\)，区域为
      \[
        -\frac{\pi}{2}\le \theta\le \frac{\pi}{2},\quad 0\le r\le 2a\cos\theta.
      \]
      又 \(z=r\)，故
      \[
        I=\sqrt2\int_{-\pi/2}^{\pi/2}\int_0^{2a\cos\theta} r^3\cos\theta\,\mathrm{d}r\,\mathrm{d}\theta
        =\frac{64\sqrt2}{15}a^4.
      \]

    \item \textbf{Solution}. 由高斯公式
      \[
        I=\iiint_V\left(\frac{\partial x^3}{\partial x}+\frac{\partial y^3}{\partial y}+\frac{\partial z}{\partial z}\right)\,\mathrm{d}v
        =\iiint_V(3x^2+3y^2+1)\,\mathrm{d}v.
      \]
      再利用球域的轮换对称性，
      \[
        \iiint_V(3x^2+3y^2)\,\mathrm{d}v=\iiint_V(x^2+y^2+z^2)\,\mathrm{d}v.
      \]
      因而
      \[
        I=\iiint_V(x^2+y^2+z^2)\,\mathrm{d}v+\iiint_V1\,\mathrm{d}v.
      \]
      用球坐标可得
      \[
        \iiint_V(x^2+y^2+z^2)\,\mathrm{d}v=\frac{4\pi R^5}{5},\qquad
        \iiint_V1\,\mathrm{d}v=\frac{4\pi R^3}{3}.
      \]
      所以
      \[
        I=\frac{4\pi R^5}{5}+\frac{4\pi R^3}{3}.
      \]

    \item \textbf{Solution}. 由比值法，收敛半径为 \(1\). 当 \(x=\pm1\) 时，通项 \(\dfrac{n}{n+1}x^n\) 不趋于 \(0\)，故两端点都发散，所以收敛域为
      \[
        (-1,1).
      \]
      对 \(|x|<1\)，
      \[
        S(x)=\sum_{n=1}^{\infty}\frac{n}{n+1}x^n
        =\sum_{n=1}^{\infty}x^n-\sum_{n=1}^{\infty}\frac{x^n}{n+1}.
      \]
      又
      \[
        \sum_{n=1}^{\infty}x^n=\frac{x}{1-x},
      \]
      以及
      \[
        \sum_{n=1}^{\infty}\frac{x^n}{n+1}
        =\frac{-\ln(1-x)-x}{x}.
      \]
      因而
      \[
        S(x)=\frac{x}{(1-x)^2}+\frac{\ln(1-x)}{x}+1\qquad (x\neq 0),
      \]
      并且 \(S(0)=0\).
  \end{enumerate}

  \section{综合题(每小题 7 分，共 14 分)}
  \begin{enumerate}
    \setcounter{enumi}{16}
    \item \textbf{Solution}. 由齐次方程通解
      \[
        y=(C_1+C_2x)e^x
      \]
      可知特征方程有二重根 \(r=1\)，于是
      \[
        y''-2y'+y=0.
      \]
      因而原方程为
      \[
        y''-2y'+y=x.
      \]
      设特解为 \(y^*=Ax+B\)，代入得
      \[
        Ax+B-2A=x,
      \]
      比较系数得 \(A=1,\ B=2\). 所以通解为
      \[
        y=(C_1+C_2x)e^x+x+2.
      \]
      由条件
      \[
        y(0)=2,\qquad y'(0)=0
      \]
      解得 \(C_1=0,\ C_2=-1\). 因而
      \[
        y=-xe^x+x+2.
      \]

    \item \textbf{Solution}. 在区域内部，
      \[
        f_x=2x(1-y^2),\qquad f_y=2y(2-x^2).
      \]
      解驻点方程得 \((0,0),(\pm\sqrt2,1)\)，其函数值分别为
      \[
        f(0,0)=0,\qquad f(\pm\sqrt2,1)=2.
      \]
      边界分两部分考察.

      当 \(y=0\) 时，
      \[
        f(x,0)=x^2,\qquad -2\le x\le 2,
      \]
      所以最小值为 \(0\)，最大值为 \(4\).

      当 \(x^2+y^2=4,\ y\ge 0\) 时，代入 \(y^2=4-x^2\) 得
      \[
        f=x^2+2(4-x^2)-x^2(4-x^2)=x^4-5x^2+8.
      \]
      比较边界与内部函数值，得全区域上
      \[
        \max f=8,\qquad \min f=0.
      \]
  \end{enumerate}

  \section{证明题(每小题 5 分，共 10 分)}
  \begin{enumerate}
    \setcounter{enumi}{18}
    \item \textbf{Proof}. 设
      \[
        a_n=\frac{(-1)^n}{n-(-1)^n}.
      \]
      先看绝对收敛性：
      \[
        |a_n|=\frac{1}{n-(-1)^n}\sim \frac{1}{n}\qquad (n\to\infty),
      \]
      故 \(\sum |a_n|\) 发散，所以原级数不绝对收敛.

      再将其写成
      \[
        a_n=\frac{(-1)^n}{n}+b_n,
      \]
      其中 \(b_n=O(n^{-2})\)，因此 \(\sum b_n\) 绝对收敛；而
      \[
        \sum \frac{(-1)^n}{n}
      \]
      是交错调和级数，条件收敛. 所以原级数条件收敛.

    \item \textbf{Proof}. 记
      \[
        P(x,y)=\frac{1}{y}[1+y^2f(xy)],\qquad
        Q(x,y)=\frac{x}{y^2}[y^2f(xy)-1].
      \]
      直接计算得
      \[
        P_y=f(xy)-\frac{1}{y^2}+xyf'(xy),\qquad
        Q_x=f(xy)-\frac{1}{y^2}+xyf'(xy),
      \]
      即 \(P_y=Q_x\). 因而在上半平面 \(y>0\) 内该积分与路径无关.

      取势函数
      \[
        \Phi(x,y)=\frac{x}{y}+F(xy),\qquad F'(t)=f(t),
      \]
      则
      \[
        \mathrm{d}\Phi=P\,\mathrm{d}x+Q\,\mathrm{d}y.
      \]
      因而
      \[
        I=\Phi(c,d)-\Phi(a,b)=\frac{c}{d}-\frac{a}{b}+F(cd)-F(ab).
      \]
      当 \(ab=cd\) 时，最后两项相消，所以
      \[
        I=\frac{c}{d}-\frac{a}{b}.
      \]
  \end{enumerate}
\end{document}






